Los experimentos se realizaron en un sistema con las siguientes especificaciones:

\subsection*{Especificaciones}
\begin{itemize}
  \item \textbf{Equipo:} \texttt{LENOVO ThinkPad X13 Gen 3}.
  \item \textbf{CPU:} \texttt{Intel Core i7-1265U (4.8 GHz)}.
  \item \textbf{RAM:} \texttt{16 GB}.
  \item \textbf{Disco:} \texttt{SSD NVMe 512 GB}.
  \item \textbf{SO:} \texttt{Ubuntu 26.04 (GNOME Edition)} con \texttt{Linux 7.0.0-22-generic}.
  \item \textbf{Compilador:} \texttt{g++ 15.2.0} con flags \texttt{-std=c++20 -O2 -g -Wall -Wextra}.
\end{itemize}

La compilación y ejecución está automatizada con \texttt{Makefile}. Los tiempos se midieron en nanosegundos $(ns)$ con \texttt{<chrono>} \parencite{cppreference:chrono} y uso de memoria en kilobytes (KB) mediante el pico máximo de memoria heap asignada por cada algoritmo. \parencite{so:memory-tracking-new-delete}